\documentclass{article}
\usepackage{graphicx} % Required for inserting images
\usepackage{booktabs}
\usepackage{enumitem}
\usepackage[table]{xcolor}
\usepackage{array}
\usepackage{natbib}
\usepackage{rotating}
\usepackage{lscape}
\usepackage{amsmath}
\usepackage{changepage}
\usepackage{pdflscape}
\usepackage{listings}
\lstset{
  language=Python,
  basicstyle=\ttfamily\small,
  keywordstyle=\color{black},
  commentstyle=\color{red},
  stringstyle=\color{red},
  showstringspaces=false,
  breaklines=true,
  captionpos=b
}

\usepackage{xcolor}

\definecolor{codegreen}{rgb}{0,0.6,0}
\definecolor{codegray}{rgb}{0.5,0.5,0.5}
\definecolor{codepurple}{rgb}{0.58,0,0.82}
\definecolor{backcolour}{rgb}{0.95,0.95,0.92}

\usepackage{longtable}
\usepackage{setspace}
\usepackage{hyperref}
\hypersetup{
    colorlinks=true,
    linkcolor=black,
    filecolor=black,
    urlcolor=black,
    citecolor=black
}

\title{Prompt Stability Scoring for Text Annotation with Large Language Models}
\author{Christopher Barrie,$^{1}$ Elli Palaiologou,$^{2}$ Petter Törnberg$^{3}$\\
\\
\normalsize{$^{1}$Department of Sociology, New York University}\\
\normalsize{$^{2}$Independent Researcher}\\
\normalsize{$^{3}$Institute for Logic, Language, and Computation, University of Amsterdam}}
\date{\today}

\begin{document}

\maketitle

\begin{abstract}
\noindent Researchers are increasingly using language models (LMs) for text annotation. These approaches rely only on a prompt telling the model to return a given output according to a set of instructions. The reproducibility of LM outputs may nonetheless be vulnerable to small changes in the prompt design. This calls into question the replicability of classification routines. To tackle this problem, researchers have typically tested a variety of semantically similar prompts to determine what we call ``prompt stability." These approaches remain ad-hoc and task specific. In this article, we propose a general framework for diagnosing prompt stability by adapting traditional approaches to intra- and inter-coder reliability scoring. We call the resulting metric the Prompt Stability Score (PSS) and provide a Python package \texttt{promptstability} for its estimation. Using six different datasets and twelve outcomes, we classify $\sim$3.1m rows of data and $\sim$300m input tokens to: a) diagnose when prompt stability is low; and b) demonstrate the functionality of the package. We conclude by providing best practice recommendations for applied researchers.

\end{abstract}

\vfill % Push the content to fill the page

\end{document}